\documentclass[11pt]{article}
\usepackage[T1]{fontenc}
\usepackage{lmodern}
\usepackage{amsmath,amssymb,amsthm,mathtools}
\usepackage{booktabs}
\usepackage[margin=1in]{geometry}
\usepackage{microtype}
\usepackage{hyperref}
\hypersetup{colorlinks=true,linkcolor=blue,citecolor=blue,urlcolor=blue,
  pdftitle={A second-order and moment proof of the five-point Mordell inequality},
  pdfauthor={Fengcan Jiang}}
\setlength{\emergencystretch}{2em}
\newtheorem{theorem}{Theorem}[section]
\newtheorem{proposition}[theorem]{Proposition}
\newtheorem{lemma}[theorem]{Lemma}
\newtheorem{corollary}[theorem]{Corollary}
\theoremstyle{remark}
\newtheorem{remark}[theorem]{Remark}
\newcommand{\C}{\mathbb C}
\newcommand{\R}{\mathbb R}
\newcommand{\Q}{\mathbb Q}
\newcommand{\one}{\mathbf 1}
\newcommand{\norm}[1]{\lVert#1\rVert}
\newcommand{\ip}[2]{\langle#1,#2\rangle}
\newcommand{\pos}[1]{[#1]_+}
\newcommand{\op}{\mathrm{op}}
\newcommand{\Rea}{\operatorname{Re}}
\newcommand{\diag}{\operatorname{diag}}
\numberwithin{equation}{section}
\title{A second-order and moment proof\\of the five-point Mordell inequality}
\author{Fengcan Jiang}
\date{October 3, 2026}

\begin{document}
\maketitle
\begin{abstract}
We give a computer-assisted proof of the sharp five-point Mordell
discriminant inequality by using second-order conditions and low-degree
moments at a maximizer. A general variational estimate bounds the variance
of the squared radii and the second complex power sum. For five points,
a two-dimensional projection supplies a constraint connecting the third
and fourth power sums. A quartic majorant for the logarithm then reduces
the global comparison to a two-variable inequality. Near the unit circle,
an angular concavity estimate and a weighted Laplacian give a separate
exclusion argument whose constants are verified algebraically.
One scalar branch is excluded by a uniform rational estimate; the
remaining branch is certified by exact rational interval arithmetic
in a short Python program using only the standard library.
The argument is independent of the classification of stationary
configurations. We also give a moment-matrix hierarchy for arbitrary
point counts and an explicit local radial gap for six points.
\end{abstract}

\noindent\textbf{Keywords:} discriminant inequality; Vandermonde product;
second-order optimality; moment bounds; exact interval arithmetic.

\noindent\textbf{2020 Mathematics Subject Classification:}
26D15, 30C10, 65G30.

\section{Introduction}
For $n\ge2$ and $z=(z_1,\ldots,z_n)\in\C^n$, write
\[
 \Delta_n(z)=\prod_{i<j}|z_i-z_j|^2,\qquad
 \mathcal S_n=\left\{z:\sum_i|z_i|^2=n\right\}.
\]
The problem considered by Mordell \cite{Mordell} asks when the regular
$n$-gon on the unit circle maximizes $\Delta_n$ on $\mathcal S_n$.
For a monic polynomial with roots $z_i$, this product is the absolute
value of its discriminant.

Wang \cite{Wang} gives a proof of the five-point inequality by exact
algebraic classification of the stationary equations, followed by
discriminant comparison. The present paper gives a different proof of
the same sharp inequality. Its starting point is the second-order
condition at a maximizer. The analytic reduction uses a complex
symmetric matrix, a two-dimensional projection, and a quartic logarithm
majorant; the final certificate involves only two real variables.

\begin{theorem}[Five-point inequality; cf.\ \cite{Wang}]
\label{thm:main}
If $z_1,\ldots,z_5\in\C$ and $\sum_i|z_i|^2=5$, then
\begin{equation}\label{eq:main}
 \Delta_5(z)\le5^5.
\end{equation}
Equality holds precisely when, after a permutation,
$z_j=e^{i\theta}e^{2\pi i(j-1)/5}$ for some $\theta\in\R$.
\end{theorem}

The proof is computer-assisted. Every reduction from
\eqref{eq:main} to the scalar inequalities is given below.
The ancillary program \texttt{verify\_mordell.py} certifies the constants
and the final scalar inequality with rational arithmetic and integer
square roots. It requires no computer algebra system or numerical
optimizer. A separate ancillary script checks the distance-moment
expansions and the polynomial encoding.

The general second-order estimates admit a hierarchy of higher-moment
constraints, developed in Section \ref{sec:extensions}. The local
angular argument also extends to six points. These extensions concern
the structure of local maximizers; the bound $\Delta_n\le n^n$ itself
fails for $n\ge6$, as shown by Chalk \cite{Chalk}.

\section{Second-order estimates}
\subsection{A general radial bound}
We use the Hermitian inner product $\ip{x}{y}=\sum_i\overline{x_i}y_i$;
a superscript $\mathsf T$ denotes transpose, and $*$ denotes conjugate
transpose. Powers of vectors below are taken coordinatewise.

\begin{proposition}\label{prop:general}
Let $n\ge2$, and let $z$ be a collision-free local maximizer of
$\Delta_n$ on $\mathcal S_n$. Set
\[
 t_i=|z_i|^2,\quad s=\sum_i(t_i-1)^2,\quad
 p_2=\sum_i z_i^2,\quad \lambda=\frac{n-1}{2}.
\]
Then
\begin{equation}\label{eq:stationary-general}
 \sum_{j\ne i}\frac1{z_i-z_j}=\lambda\overline{z_i},
 \qquad \sum_i z_i=0.
\end{equation}
The complex symmetric matrix
\begin{equation}\label{eq:B}
 B_{ij}=-\frac1{(z_i-z_j)^2}\quad(i\ne j),\qquad
 B_{ii}=\sum_{j\ne i}\frac1{(z_i-z_j)^2}
\end{equation}
satisfies
\begin{equation}\label{eq:B-general}
 B\one=0,\qquad Bz=\lambda\overline z,\qquad
 \norm{B}_{\op}\le\lambda.
\end{equation}
Moreover,
\begin{equation}\label{eq:radial-general}
 3ns+|p_2|^2\le n^2,\qquad
 |z_i-z_j|^2\ge\frac{2n}{(n-1)^2}.
\end{equation}
If $n\ge3$, the second power sum also satisfies
\begin{equation}\label{eq:power-radial}
 |p_2|^2\le
 \frac{ns(n+s)}{s+n(n-2)^2/(n-1)^2}.
\end{equation}
For the real weighted Laplacian with weights
$|z_i-z_j|^{-2}$, one has
\begin{equation}\label{eq:L-general}
 Lt=(n-1)(t-\one).
\end{equation}
\end{proposition}

\begin{proof}
Lagrange multipliers for $\log\Delta_n$ give the first equation
in \eqref{eq:stationary-general} with a real multiplier.
Multiplying by $z_i$ and summing yields
$\lambda n=n(n-1)/2$; summing the equations gives the zero centroid.
The first two identities in \eqref{eq:B-general} follow directly.

For $\ip{z}{h}=0$, every phase multiple of $h$ is a real tangent vector
to $\mathcal S_n$. The second derivative of the Lagrangian is
\[
 -2\Rea(h^{\mathsf T}Bh)-2\lambda\norm{h}^2.
\]
Its nonpositivity for every phase implies
$|h^{\mathsf T}Bh|\le\lambda\norm h^2$.
Write any vector as $h+\alpha z$ with $\ip z h=0$.
Since $h^{\mathsf T}Bz=0$ and $z^{\mathsf T}Bz=\lambda n$,
the same quadratic bound holds for all vectors.
For unit vectors $x,y$, the identity
\[
4x^{\mathsf T}By=(x+y)^{\mathsf T}B(x+y)
 -(x-y)^{\mathsf T}B(x-y)
\]
and the parallelogram law give $|x^{\mathsf T}By|\le\lambda$.
Taking the supremum proves the operator norm bound.

Pairwise expansion gives
$Bz^{\circ2}=(n-1)(t-\one)$. Applying \eqref{eq:B-general} to
$z^{\circ2}-p_2\one/n$ yields
\[
 (n-1)^2s\le\lambda^2
 \left(\sum_i t_i^2-\frac{|p_2|^2}{n}\right).
\]
As $\sum_i t_i^2=n+s$, this is the first bound in
\eqref{eq:radial-general}. To sharpen the entry bound, let
$P_0=I-\one\one^{\mathsf T}/n$. Since $B=P_0BP_0$ and
$\norm{P_0e_i}^2=(n-1)/n$, we have
$|B_{ij}|\le\lambda(n-1)/n$.
This proves the second inequality.
To prove \eqref{eq:power-radial}, multiply the stationary equations
by $z_i^3$ and sum. Since $\sum_i z_i=0$, pairwise expansion gives
$\sum_i t_i z_i^2=(2n-3)p_2/(n-1)$. Thus
\[
 \ip{t-\one}{z^{\circ2}-p_2\one/n}
   =\frac{n-2}{n-1}p_2.
\]
Cauchy--Schwarz yields
$((n-2)/(n-1))^2|p_2|^2
 \le s(n+s-|p_2|^2/n)$.
Rearranging proves \eqref{eq:power-radial}, including $s=0$.
Finally, multiply \eqref{eq:stationary-general} by $z_i$, take real
parts, and use
$2\Rea(z_i/(z_i-z_j))=1+(t_i-t_j)/|z_i-z_j|^2$.
This gives \eqref{eq:L-general}.
\end{proof}

\subsection{Five-point moments and a Gram matrix}
From now on $n=5$, and we work at a collision-free local maximizer.
In addition to $t_i$ and $s$, put
\begin{equation}\label{eq:notation}
\begin{gathered}
 T=\sum_i t_i^2=5+s,\qquad U=\sum_i t_i^3,\qquad V=\sum_i t_i^4,\\
 p_j=\sum_i z_i^j,\quad q=|p_2|^2,\quad p=|p_3|^2,\quad
 w=|p_4|,\quad m=\sum_i t_i^2z_i^2.
\end{gathered}
\end{equation}
Proposition \ref{prop:general} gives
\begin{equation}\label{eq:domain}
 0\le s\le\frac53,\qquad
 0\le q\le \min\left\{25-15s,\frac{80s(5+s)}{45+16s}\right\}.
\end{equation}
Multiplication of the stationary equations by $z_i^2,z_i^3,z_i^4$,
followed by summation and pairing, gives
\begin{equation}\label{eq:mixed}
 \sum_i t_i z_i=0,\qquad
 \sum_i t_i z_i^2=\frac74p_2,\qquad
 \sum_i t_i z_i^3=\frac32p_3.
\end{equation}
For example, for the middle identity the pair sum is
$4p_2+\sum_{i<j}z_i z_j=7p_2/2$.
Direct expansion also gives
\begin{equation}\label{eq:B-powers}
 Bz^{\circ2}=4(t-\one),\qquad
 Bz^{\circ4}=(8t-9\one)\circ z^{\circ2}-p_2\one.
\end{equation}

The matrix $4I-B^*B$ is positive semidefinite.
Let $D$ be its Gram matrix on the two vectors
$z^{\circ2}-p_2\one/5$ and $z^{\circ4}-p_4\one/5$.
Using \eqref{eq:mixed}--\eqref{eq:B-powers}, its entries are
\begin{equation}\label{eq:D}
\begin{aligned}
 D_{22}&=20-12s-\frac45q,\\
 D_{44}&=-60V+144U-81T+5q-\frac45w^2,\\
 D_{24}&=-28m+83p_2-\frac45\overline{p_2}p_4.
\end{aligned}
\end{equation}
In particular,
\begin{equation}\label{eq:D-psd}
 D_{22},D_{44}\ge0,\qquad |D_{24}|^2\le D_{22}D_{44}.
\end{equation}

\section{A two-dimensional projection constraint}
\begin{lemma}\label{lem:projection}
For $s>0$ at a collision-free local maximizer,
\begin{equation}\label{eq:w-lower}
 w\ge T-\frac{2q}{5}-\frac{9q}{8s}-\frac{p}{5-\sqrt q}.
\end{equation}
\end{lemma}
\begin{proof}
By \eqref{eq:domain}, $q<25$.
Thus $E=\operatorname{span}_{\C}\{\one,z,\overline z\}$ has
dimension three: $\one$ is orthogonal to the other two vectors, whose
Gram determinant is $25-q>0$.
The vector $u=t-\one$ is real, lies in $E^\perp$, and has norm squared $s$.
Let $v$ be the orthogonal projection of $z^{\circ2}$ onto $E^\perp$.
The Gram matrix of $z,\overline z$ is
$\left(\begin{smallmatrix}5&\overline{p_2}\\p_2&5\end{smallmatrix}\right)$.
Since $\ip z {z^{\circ2}}=0$ and
$\ip{\overline z}{z^{\circ2}}=p_3$, the linear part of the
projection onto $E$ is
\[
 \frac{-\overline{p_2}p_3z+5p_3\overline z}{25-q}.
\]
Consequently,
\begin{equation}\label{eq:projection-identities}
\begin{aligned}
 \norm v^2&=T-\frac q5-\frac{5p}{25-q},&
 \ip u v&=\frac34p_2,\\
 v^{\mathsf T}v&=p_4-\frac{p_2^2}{5}
    +\frac{\overline{p_2}p_3^2}{25-q}.
\end{aligned}
\end{equation}
The conjugation-invariant space $E^\perp$ has a real orthonormal basis
$u/\sqrt s,e$. Write $v=a\,u/\sqrt s+b e$.
Then $|a|^2=9q/(16s)$ and
\[
 |v^{\mathsf T}v|=|a^2+b^2|
 \ge\norm v^2-2|a|^2
 =\norm v^2-\frac{9q}{8s}.
\]
Combine this with \eqref{eq:projection-identities} and the triangle
inequality. The coefficient of $p$ becomes
$(5+\sqrt q)/(25-q)=1/(5-\sqrt q)$, proving the assertion.
\end{proof}

\section{Excluding small nonzero radial variance}
\begin{proposition}[Local radial gap]\label{prop:gap}
A collision-free local maximizer satisfies either $s=0$ or $s>1/20$.
\end{proposition}
\begin{proof}
Suppose $0<s\le1/20$, and write $z_i=r_i e^{i\theta_i}$.
For $\delta_i=t_i-1$, we have $\sum_i\delta_i=0$ and
$\sum_i\delta_i^2=s$. Hence
\begin{equation}\label{eq:local-radii}
 t_i\ge1-\sqrt{4s/5}\ge\frac45,\qquad
 \frac78<r_ir_j<\frac98.
\end{equation}
For completeness, set $a=(t_i+t_j)/2$, $b=(t_i-t_j)/2$.
Cauchy--Schwarz for the remaining three deviations gives
$s\ge(10/3)(a-1)^2+2b^2$.
It follows that
$t_i t_j\ge a^2-s/2+(5/3)(a-1)^2$.
This last expression is increasing for
$a\ge1-\sqrt{3s/10}>7/8$, since its derivative is $(16a-10)/3$.
At the left endpoint it equals $a^2$, so
$r_ir_j\ge1-\sqrt{3s/10}>7/8$.
The upper bound follows from $r_ir_j\le a\le1+\sqrt{3s/10}<9/8$;
here $(1/8)^2>3/200$.

Let $\overline r=5^{-1}\sum_i r_i$ and
$\sigma=\sum_i(r_i-\overline r)^2$.
Using $(r_i+r_j)^2\ge4r_ir_j>7/2$,
\begin{equation}\label{eq:sigma}
 \sigma=\frac15\sum_{i<j}
       \frac{(t_i-t_j)^2}{(r_i+r_j)^2}
 <\frac{2s}{7}\le\frac1{70},\qquad
 (r_i-r_j)^2<\frac{4s}{7}\le\frac1{35}.
\end{equation}
All squared distances are at least $1/2$ by
Proposition \ref{prop:general}. Thus
\[
 \sin^2\frac{\theta_i-\theta_j}{2}
 =\frac{|z_i-z_j|^2-(r_i-r_j)^2}{4r_ir_j}
 >\frac{11}{105}>\frac1{10}.
\]
The angles are distinct and all the corresponding half-angle cotangents
have absolute value less than $3$.

Lift and order the angles so that
$\theta_1<\cdots<\theta_5<\theta_1+2\pi$, and put
\[
 \mathcal C(\theta)=\sum_{i<j}
       \log|e^{i\theta_i}-e^{i\theta_j}|^2.
\]
For $\sum_i h_i=0$,
\[
 h^{\mathsf T}\nabla^2\mathcal C(\theta)h
 =-\frac12\sum_{i<j}
   \csc^2\frac{\theta_i-\theta_j}{2}(h_i-h_j)^2
 \le-\frac52\norm h^2.
\]
Let $\theta^*$ be the equally spaced angles with the same mean, and let
$\eta=\norm{\theta-\theta^*}$. The ordered region of span less than
$2\pi$, intersected with the fixed-mean hyperplane, is convex.
Integration of the preceding Hessian bound along the connecting segment
gives
\begin{equation}\label{eq:eta-gradient}
 \eta\le\frac25\norm{\nabla\mathcal C(\theta)}.
\end{equation}
The angular stationary equations for the actual radii imply
\begin{equation}\label{eq:angular-stationary}
 \frac{\partial\mathcal C}{\partial\theta_i}
 =\sum_{j\ne i}\cot\frac{\theta_i-\theta_j}{2}
       \frac{(r_i-r_j)^2}{|z_i-z_j|^2}.
\end{equation}
The incidence matrix of $K_5$ has operator norm $\sqrt5$.
For $v_i=r_i-\overline r$, the zero sum gives
$\max_i v_i^2\le4\sigma/5$, and expansion yields
$\sum_{i<j}(v_i-v_j)^4=5\sum_i v_i^4+3\sigma^2\le7\sigma^2$.
Combining these observations with
\eqref{eq:eta-gradient}--\eqref{eq:angular-stationary}, if squared
distances are at least $d$ and cotangents at most $h$ in absolute value,
then
\begin{equation}\label{eq:eta-basic}
 \eta\le\gamma\sigma h/d,\qquad \gamma=\frac{2\sqrt{35}}5.
\end{equation}
Initially $d=1/2,h=3$ and \eqref{eq:sigma} give
$\eta<3\gamma/35<21/100$.

We next use only algebraic constants. Put
\begin{equation}\label{eq:algebraic-trig}
\begin{gathered}
 S_1=\sqrt{\frac{5-\sqrt5}{8}},\quad C_1=\frac{1+\sqrt5}{4},
 \qquad
 S_2=\sqrt{\frac{5+\sqrt5}{8}},\quad C_2=\frac{\sqrt5-1}{4},\\
 P_j(u)=S_j(1-u^2/4)-C_j u/\sqrt2,\qquad
 N(u)=C_1+S_1u/\sqrt2.
\end{gathered}
\end{equation}
Here $S_j=\sin(j\pi/5)$, $C_j=\cos(j\pi/5)$.
The inequalities $\cos x\ge1-x^2/2$ and $\sin x\le x$ for $x\ge0$
give
\[
 \sin(j\pi/5-u/\sqrt2)\ge P_j(u),\qquad
 \cos(\pi/5-u/\sqrt2)\le N(u).
\]
If $\eta\le u\le21/100$, each pairwise angular error is at most
$\sqrt2u$. Consequently every absolute half-angle sine is at least
$P_1(u)>0$ and the absolute cotangent is at most $N(u)/P_1(u)$.
Equations \eqref{eq:local-radii}, \eqref{eq:sigma}, and
\eqref{eq:eta-basic} then give
\begin{equation}\label{eq:algebraic-iteration}
 \eta<\frac{\gamma N(u)}{245P_1(u)^3}.
\end{equation}
The exact algebraic checks in the ancillary program prove
\begin{equation}\label{eq:two-steps}
 \frac{\gamma N(21/100)}{245P_1(21/100)^3}<\frac9{100},
 \qquad
 \frac{\gamma N(9/100)}{245P_1(9/100)^3}<\frac3{50}.
\end{equation}
Thus $\eta<3/50$. The same algebraic checks give
\begin{equation}\label{eq:weights}
 \frac{2}{7P_1(3/50)^2}<\frac{19}{20},\qquad
 \frac{2}{7P_2(3/50)^2}<\frac13.
\end{equation}
These bound the inverse squared distances on consecutive pairs
and diagonals in the cyclic order, respectively.

Let $L$ be the real Laplacian with weights $|z_i-z_j|^{-2}$.
Comparison in the positive semidefinite order with a pentagon of
edge weight $19/20$ and diagonal weight $1/3$ gives
\begin{equation}\label{eq:spectral-gap}
 \lambda_{\max}(L)\le
 \frac{19}{20}\frac{5+\sqrt5}{2}
 +\frac13\frac{5-\sqrt5}{2}<4.
\end{equation}
The two Laplacians of the pentagon and its complement are simultaneously
diagonalized by the discrete Fourier basis; their nonzero eigenvalues
are $(5\pm\sqrt5)/2$ in complementary order.
However, \eqref{eq:L-general} gives $L(t-\one)=4(t-\one)$.
As $s>0$, this contradicts \eqref{eq:spectral-gap}.
\end{proof}

\section{A quartic logarithm majorant}
Put
\begin{equation}\label{eq:majorant-constants}
 a=\frac{5-\sqrt5}{2},\quad b=\frac{5+\sqrt5}{2},\quad
 c=\frac15-\frac{2}{5\sqrt5}\log\frac ba,\quad k=\frac9{2500}.
\end{equation}
The following bounds are certified in the ancillary program:
\begin{equation}\label{eq:c-bounds}
 \frac{22}{3}k<c<C:=\frac{278365}{10000000}<\frac{17}{2}k.
\end{equation}
Let $H$ be the cubic interpolating $\log x$ and its derivative at $a,b$.
Solving the four interpolation conditions gives
\begin{equation}\label{eq:H}
 H(x)=\frac12\log5-\frac74-\frac{25c}{4}
 +(1+15c)x-\left(\frac1{10}+\frac{15c}{2}\right)x^2+cx^3.
\end{equation}

\begin{lemma}\label{lem:majorant}
For $0<x\le R:=41/6$,
\begin{equation}\label{eq:majorant}
 \log x\le H(x)-k(x^2-5x+5)^2.
\end{equation}
\end{lemma}
\begin{proof}
Add the node $R$ to the double nodes $a,b$.
The quartic Hermite interpolant is
$H(x)-k_R(x^2-5x+5)^2$, where
\[
 k_R=\frac{H(R)-\log R}{(R^2-5R+5)^2}.
\]
The fifth derivative of $\log x$ is positive.
Its Hermite remainder has the sign of
$(x-a)^2(x-b)^2(x-R)$, hence is nonpositive for $0<x\le R$.
The exact series check in Appendix \ref{app:certificate} proves
$k_R>k$. This yields \eqref{eq:majorant}.
\end{proof}

All squared distances at a local maximizer belong to this range:
\[
 |z_i-z_j|^2\le2(t_i+t_j)
 \le4+2\sqrt{6s/5}\le4+2\sqrt2<41/6.
\]
Write $d_{ij}=|z_i-z_j|^2$ and
$\mathcal Q=\sum_{i<j}(d_{ij}^2-5d_{ij}+5)^2$.
Expansion using \eqref{eq:mixed} gives
\begin{equation}\label{eq:distance-moments}
\begin{aligned}
 \sum_{i<j}d_{ij}&=25,\\
 \sum_{i<j}d_{ij}^2&=5T+50+q,\\
 \sum_{i<j}d_{ij}^3&=5U+45T+\frac{21}{2}q-p,\\
 \sum_{i<j}d_{ij}^4&=5V+80U+18T^2+49q-12p+w^2
                 +12\Rea(m\overline{p_2}).
\end{aligned}
\end{equation}
These identities can be checked uniformly by putting
$M_{\alpha\beta}=\sum_i z_i^\alpha\overline{z_i}^{\,\beta}$ and using
\begin{equation}\label{eq:moment-expansion}
 \sum_{i<j}d_{ij}^r=
 \frac12\sum_{\alpha,\beta=0}^r(-1)^{\alpha+\beta}
 \binom r\alpha\binom r\beta
 M_{r-\alpha,r-\beta}M_{\alpha\beta}.
\end{equation}

The regular pentagon has five squared distances $a$ and five $b$.
Subtract its value in \eqref{eq:majorant} and use
\eqref{eq:distance-moments} to obtain
\begin{equation}\label{eq:log-bound}
 \log\Delta_5-5\log5\le
 \left(-\frac12+\frac{15c}{2}\right)s+5c(U-5)
 +\left(-\frac1{10}+3c\right)q-cp-k\mathcal Q.
\end{equation}
The moment expansion of $\mathcal Q$ is
\begin{equation}\label{eq:Q-expand}
\begin{aligned}
 \mathcal Q={}&5V+30U+18T^2-275T+750-21q-2p+w^2\\
             &+12\Rea(m\overline{p_2}).
\end{aligned}
\end{equation}
Set
\[
 R_0=5V+30U+18T^2-275T+750,\qquad
 D_0=-60V+144U-81T+5q.
\]
Solving \eqref{eq:D} for $m$ and applying \eqref{eq:D-psd} gives
\begin{equation}\label{eq:Q-lower}
 \mathcal Q\ge R_0+\frac{102}{7}q-2p+w^2-\frac{12}{35}qw
       -\frac37\sqrt{qD_{22}D_0}.
\end{equation}
Indeed, the contribution of $m$ is
\[
12\Rea(m\overline{p_2})
 =\frac{249}{7}q-\frac{12}{35}\Rea(\overline{p_2}^{\,2}p_4)
 -\frac37\Rea(D_{24}\overline{p_2}).
\]
Also $D_0=D_{44}+4w^2/5\ge D_{44}\ge0$, so replacing $D_{44}$
by $D_0$ in the square root weakens the lower bound in the required
direction.

\section{Reduction to two variables}
\subsection{Eliminating the complex power sums}
Assume $s>0$. Put $L_0=T-2q/5-9q/(8s)$ and $w_0=6q/35$.
Lemma \ref{lem:projection} implies
$p\ge(5-\sqrt q)\pos{L_0-w}$.
The expression to be bounded below after combining
\eqref{eq:log-bound} and \eqref{eq:Q-lower} is
$k(w^2-12qw/35)+(c-2k)p$.
It satisfies
\begin{equation}\label{eq:eliminate-pw}
 k(w^2-12qw/35)+(c-2k)p
 \ge k\left(\pos{T-\left(\frac47+\frac9{8s}\right)q}^{\,2}
                  -\frac{36q^2}{1225}\right).
\end{equation}
To see this, complete the square at $w_0$.
If $L_0\le w_0$, the unconstrained minimum at $w_0$ is feasible.
If $L_0>w_0$, the minimum is at $L_0$ provided
\[
 L_0-w_0\le\frac{c-2k}{2k}(5-\sqrt q).
\]
Here $L_0>0$ implies
$q<T/(2/5+9/(8s))\le800/129<25/4$.
By \eqref{eq:c-bounds}, the last right-hand side is greater than
$(8/3)(5/2)=20/3\ge T\ge L_0-w_0$.
This proves the asserted piecewise minimum.

\subsection{Eliminating the radial moments}
Cauchy--Schwarz and projection of $t^{\circ2}$ onto the orthogonal
vectors $\one,t-\one$ give
\begin{equation}\label{eq:UV}
 U\ge\frac{T^2}{5},\qquad
 V\ge V_0(U,s):=\frac{T^2}{5}+\frac{(U-T)^2}{s}.
\end{equation}
The zero-sum deviations $\delta_i=t_i-1$ also satisfy
\[
 \sum_i\delta_i^3\ge-\frac3{\sqrt{20}}s^{3/2}.
\]
For a proof, at an extremum of the third moment on
$\sum\delta_i=0$, $\sum\delta_i^2=s$, the multiplier equation is a
quadratic in $\delta_i$. Thus at most two values occur; checking
multiplicities $1+4$ and $2+3$ gives the stated sharp lower bound.
Consequently,
\begin{equation}\label{eq:U0}
 U\ge U_0(s)=
 \begin{cases}
 5+3s-\dfrac3{\sqrt{20}}s^{3/2},&0<s\le5/4,\\[2pt]
 (5+s)^2/5,&5/4\le s\le5/3.
 \end{cases}
\end{equation}
The two candidates meet at $s=5/4$, as follows by setting
$x=\sqrt{s/20}$ and factoring
$1-3x-4x^2=(1-4x)(1+x)$.

After \eqref{eq:eliminate-pw}, the right-hand side of the logarithm
bound decreases with $V$: its nonradical derivative is $-5k$ and
the radicand decreases with $V$.
We may therefore set $V=V_0(U,s)$.
For $U\ge U_0(s)$,
\[
 \partial_U V_0=\frac{2(U-T)}s
 \ge\begin{cases}
 4-3\sqrt{s/5}\ge5/2,&s\le5/4,\\
 2(5+s)/5\ge5/2,&s\ge5/4.
 \end{cases}
\]
Hence $\partial_U D_0\le-6$, while the remaining derivative is at most
$5c-(85/2)k<0$. We may also set $U=U_0(s)$.
Both downward substitutions increase $D_0$, preserving its
nonnegativity along the substitution path.
Finally the resulting upper bound increases with $c$, with derivative
$(15/2)s+5(U_0-5)+3q>0$, so we may replace $c$ by the rational upper
bound $C$ in \eqref{eq:c-bounds}.

Define
\begin{equation}\label{eq:scalar-definitions}
\begin{gathered}
 U_*=U_0(s),\quad V_*=\frac{(5+s)^2}{5}
                  +\frac{(U_*-5-s)^2}{s},\\
 A_q=\frac1{10}-3C+\frac{102k}{7},\quad
 B_q=\frac{36k}{1225},\quad D_q=\frac{3k}{7},\\
 f(s)=\left(-\frac12+\frac{15C}{2}\right)s+5C(U_*-5)\\
 \hspace{15mm}-k\{5V_*+30U_*+18(5+s)^2-275(5+s)+750\},\\
 D_*(s)=-60V_*+144U_*-81(5+s).
\end{gathered}
\end{equation}
We have proved the following scalar reduction.

\begin{proposition}\label{prop:scalar}
At a collision-free local maximizer with $s>0$,
\begin{equation}\label{eq:F-bound}
 \log\Delta_5-5\log5\le F(s,q),
\end{equation}
where
\begin{equation}\label{eq:F}
\begin{aligned}
 F(s,q)={}&f(s)-A_q q+B_q q^2\\
 &+D_q\sqrt{q(20-12s-4q/5)(D_*(s)+5q)}\\
 &-k\pos{5+s-\left(\frac47+\frac9{8s}\right)q}^{\,2}.
\end{aligned}
\end{equation}
The parameters satisfy \eqref{eq:domain}, and both nontrivial factors
inside the square root are nonnegative.
\end{proposition}

\section{Exact certificate and completion of the proof}
\begin{lemma}[Scalar certificate]\label{lem:certificate}
With the constants in \eqref{eq:scalar-definitions},
\begin{equation}\label{eq:certificate}
 F(s,q)<0\qquad
 \left(\frac1{20}\le s\le\frac53,\quad
 0\le q\le\min\left\{25-15s,\frac{80s(5+s)}{45+16s}\right\}\right).
\end{equation}
\end{lemma}

\begin{proof}
We treat \eqref{eq:certificate} in two regions:
\[
 \mathrm A:\ s=20x^2,\quad\frac1{20}\le x\le\frac14;
 \qquad
 \mathrm B:\ \frac54\le s\le\frac53.
\]
In A, $U_*=5+60x^2-60x^3$ and
$V_*=5+120x^2-240x^3+260x^4$.
In B, $U_*=(5+s)^2/5$ and $V_*=(5+s)^3/25$.
The endpoint $20(1/20)^2=1/20$ exactly matches the local gap,
and the regions meet at $s=5/4$.

Region A starts with its containing rectangle and $0\le q\le25$.
For a rectangle $[x_l,x_h]\times[q_l,q_h]$, put
$s_{\min}=20x_l^2$, $s_{\max}=20x_h^2$, and replace $q_h$ by
\[
 \min\left\{q_h,25-15s_{\min},
            \frac{80s_{\max}(5+s_{\max})}{45+16s_{\max}}\right\}.
\]
The last expression is justified by \eqref{eq:power-radial} and the
monotonicity of $80s(5+s)/(45+16s)$ for $s\ge0$.
Discard an empty rectangle. Evaluate the interval form of
\eqref{eq:F}, clipping its radicand factors to their nonnegative parts.
If the upper endpoint is negative, accept the rectangle.
Otherwise bisect its longer normalized side, using scales $1/5$ and
$1$. The depth-first search exhausts all rectangles.

\begin{center}
\begin{tabular}{lrrr}
\toprule
Region & Visited nodes & Accepted rectangles & Maximum depth\\
\midrule
A & 183 & 92 & 21\\
\bottomrule
\end{tabular}
\end{center}
No rectangles remain unresolved. Every accepted upper endpoint is less
than the following convenient rational bound:
\begin{equation}\label{eq:certificate-margins}
 -\frac1{200000}.
\end{equation}
All comparisons use exact fractions. The interval arithmetic and the
constant checks are detailed in Appendix \ref{app:certificate}.
Pruning removes only infeasible points and each bisection covers its
parent, so the finite certified cover proves the assertion.

On the second branch no subdivision is necessary. In fact,
\begin{equation}\label{eq:branch-b-analytic}
 F(s,q)\le-\frac{56657645489}{1713169664000}<-\frac3{100}
 \qquad\left(\frac54\le s\le\frac53,\quad 0\le q\le25-15s\right).
\end{equation}
To see this, write
\[
 f(s)=a_0s+b_0s^2-\frac{k}{5}s^3,\quad
 a_0=-\frac12+\frac{35}{2}C+20k,\quad b_0=C-27k.
\]
The rational inequalities $b_0<0$ and $a_0+(5/2)b_0<0$
imply that $f$ is decreasing for $s\ge5/4$. Likewise,
\[
 D_*(s)=15+27s-\frac{36}{5}s^2-\frac{12}{5}s^3
\]
is decreasing on this interval, since its derivative is at most $-9/4$.
Consequently $f(s)\le-9181/256000$, $D_*(s)\le525/16$,
$q\le25/4$, $20-12s-(4/5)q\le5$, and
$D_*(s)+5q\le1025/16$. Put
\[
 \alpha=A_q-\frac{25}{4}B_q=\frac{6692069}{98000000}>0.
\]
Dropping the nonpositive cutoff term in $F$ and completing the square
in $\sqrt q$ gives
\[
 F(s,q)\le-\frac{9181}{256000}-\alpha q
       +D_q\sqrt{\frac{5125}{16}q}
 \le-\frac{9181}{256000}
       +\frac{5125D_q^2}{64\alpha}
 =-\frac{56657645489}{1713169664000}.
\]
Every constant comparison in this argument is checked exactly in the
ancillary program.
\end{proof}

\begin{proof}[Proof of Theorem \ref{thm:main}]
The constraint sphere is compact, and a regular pentagon gives
$\Delta_5=5^5>0$. A global maximizer exists and is collision-free.
If $0<s\le1/20$, Proposition \ref{prop:gap} excludes it.
If $s\ge1/20$, Proposition \ref{prop:scalar} and
Lemma \ref{lem:certificate} give $\Delta_5<5^5$, a contradiction.
Thus $s=0$ at every global maximizer, so $|z_i|=1$.

For completeness, the classical unit-circle conclusion can be read
directly from the root polynomial, as in \cite{Mordell}.
Let $P(\zeta)=\prod_i(\zeta-z_i)$.
The stationary equations give
\[
 \frac{P''(z_i)}{2P'(z_i)}
 =\sum_{j\ne i}\frac1{z_i-z_j}=2\overline{z_i}=\frac2{z_i}.
\]
The polynomial $\zeta P''(\zeta)-4P'(\zeta)$ has degree at most three
and vanishes at five distinct points, so it is zero.
Coefficient comparison gives $P(\zeta)=\zeta^5-\alpha$, $|\alpha|=1$.
The roots are precisely a unit regular pentagon, and
$\Delta_5=\prod_i|P'(z_i)|=5^5$.
This establishes both the inequality and all equality cases.
\end{proof}

\begin{remark}
The local gap is a statement about every local maximizer, whereas
\eqref{eq:F-bound} need not exclude a local maximizer whose value is
strictly below $5^5$. The global comparison in the final proof uses the
known regular-pentagon competitor.
\end{remark}

\section{A general moment hierarchy and its limitations}
\label{sec:extensions}
The second-order method extends to arbitrary $n$, although the target
$\Delta_n\le n^n$ does not: Chalk's counterexample already disproves this
bound for every $n\ge6$ \cite{Chalk}. We record constraints that remain
valid at all collision-free local maximizers and a quantitative six-point
consequence. The polygonal spectral calculation below records a
limitation of the method; it is not presented as a new stability result.

\subsection{Higher-degree second-order constraints}
Work at a collision-free local maximizer on $\mathcal S_n$, and put
$p_k=\sum_i z_i^k$, $p_0=n$, and $\lambda=(n-1)/2$.
For every $k\ge1$, the stationary equations imply
\begin{equation}\label{eq:general-mixed-hierarchy}
 \sum_i t_i z_i^k=
 \frac{(2n-k-1)p_k+\sum_{j=1}^{k-1}p_jp_{k-j}}{n-1}.
\end{equation}
Indeed, multiply the stationary equation by $z_i^{k+1}$, sum, and pair
indices; the result before division by $\lambda$ is
$\frac12\{\sum_{j=0}^k p_jp_{k-j}-(k+1)p_k\}$.
For $k\ge2$ a related coordinatewise identity is
\begin{equation}\label{eq:general-B-hierarchy}
 Bz^{\circ k}=
 \left(k\lambda t-(k-1)(n-k/2)\one\right)\circ z^{\circ(k-2)}
 -\sum_{j=1}^{k-2}(k-1-j)p_jz^{\circ(k-2-j)}.
\end{equation}
To verify it, expand
$z_i^k-z_j^k=(z_i-z_j)\sum_{r=0}^{k-1}z_i^{k-1-r}z_j^r$
and then use, for $r\ge1$,
\[
 \sum_{j\ne i}\frac{z_j^r}{z_i-z_j}
 =\lambda z_i^r\overline z_i
 -\sum_{\ell=0}^{r-1}p_\ell z_i^{r-1-\ell}+r z_i^{r-1}.
\]

\begin{proposition}[Moment Gram hierarchy]\label{prop:moment-hierarchy}
For any finite set of positive integers $\mathcal K$, set
$x_k=z^{\circ k}-p_k\one/n$ and define
\begin{equation}\label{eq:general-D-hierarchy}
 D_{k\ell}=\lambda^2\ip{x_k}{x_\ell}
       -\ip{Bz^{\circ k}}{Bz^{\circ\ell}},\qquad k,\ell\in\mathcal K.
\end{equation}
Then $D$ is Hermitian positive semidefinite. Its entries are polynomials
in mixed moments, with the identities
\eqref{eq:general-mixed-hierarchy}--\eqref{eq:general-B-hierarchy}
providing exact reductions.
\end{proposition}
\begin{proof}
The matrix $\lambda^2I-B^*B$ is positive semidefinite by
Proposition~\ref{prop:general}, and $B\one=0$. Thus $D$ is its Gram
matrix on the vectors $x_k$. Constant projection removes a nonnegative
term that would otherwise weaken the constraints.
\end{proof}

For an explicit reusable instance, write
\[
 T=\sum_i t_i^2=n+s,\quad U=\sum_i t_i^3,\quad V=\sum_i t_i^4,
 \quad q=|p_2|^2,\quad w=|p_4|,\quad m=\sum_i t_i^2z_i^2.
\]
The submatrix on $\mathcal K=\{2,4\}$ has entries
\begin{equation}\label{eq:general-D24}
\begin{aligned}
 D_{22}&=\frac{(n-1)^2}{4}\left(n-3s-\frac qn\right),\\
 D_{44}&=-\frac{15(n-1)^2}{4}V+12(n-1)(n-2)U
       -9(n-2)^2T+nq-\frac{(n-1)^2}{4n}w^2,\\
 D_{24}&=-\frac{7(n-1)^2}{4}m+(7n^2-22n+18)p_2
       -\frac{(n-1)^2}{4n}\overline{p_2}p_4.
\end{aligned}
\end{equation}
They satisfy $D_{22},D_{44}\ge0$ and
$|D_{24}|^2\le D_{22}D_{44}$.
These formulas follow from
$Bz^{\circ4}=(2(n-1)t-3(n-2)\one)\circ z^{\circ2}-p_2\one$
and \eqref{eq:general-mixed-hierarchy}, and reduce exactly to
\eqref{eq:D} when $n=5$. Increasing $\mathcal K$ supplies additional
necessary constraints; it does not by itself prove a global upper bound.

\subsection{The dimensional obstruction in the projection step}
Assume $n\ge3$ and $s>0$, set $p=|p_3|^2$, and let
$E=\operatorname{span}_{\C}\{\one,z,\overline z\}$.
The radial bound gives $q<n^2$, so $\dim E=3$.
Let $u=t-\one$ and let $v$ be the projection of $z^{\circ2}$ onto
$E^\perp$. The computation in Lemma~\ref{lem:projection} now gives
\begin{equation}\label{eq:general-projection-identities}
\begin{aligned}
 \norm v^2&=T-\frac qn-\frac{np}{n^2-q},&
 \ip u v&=\frac{n-2}{n-1}p_2,\\
 v^{\mathsf T}v&=p_4-\frac{p_2^2}{n}
                  +\frac{\overline{p_2}p_3^2}{n^2-q}.
\end{aligned}
\end{equation}
Since $u\in E^\perp$ and $\norm u^2=s$, Cauchy--Schwarz proves the
dimension-independent bound
\begin{equation}\label{eq:general-projection-CS}
 \frac{(n-2)^2}{(n-1)^2}q
 \le s\left(T-\frac qn-\frac{np}{n^2-q}\right).
\end{equation}
Dropping the last nonpositive term recovers \eqref{eq:power-radial}.
For $n=5$, $E^\perp$ admits an orthonormal basis of two real vectors,
and the further
bound on $|v^{\mathsf T}v|$ used earlier is available.
For $n\ge6$ that particular argument fails: after fixing the real
direction $u$, two real orthonormal directions $e_1,e_2$ remain, and
$b(e_1+i e_2)$ has positive Hermitian norm but zero transpose square.
Thus no analogous lower bound follows solely from the norm and the
component along $u$. Additional moment constraints would be needed.

\subsection{A quantitative six-point radial gap}
\begin{proposition}\label{prop:six-gap}
At a collision-free local maximizer of $\Delta_6$ on $\mathcal S_6$,
either $s=0$ or $s>1/200$.
\end{proposition}
\begin{proof}
Suppose $0<s\le1/200$, and write $z_i=r_i e^{i\theta_i}$.
For $a=(t_i+t_j)/2$ and $b=(t_i-t_j)/2$, the zero-sum deviations give
$s\ge3(a-1)^2+2b^2$.
The function $a^2-s/2+\frac32(a-1)^2$ is increasing for $a>3/5$,
so its value at $a=1-\sqrt{s/3}$ proves
\[
 \frac{23}{24}<1-\sqrt{s/3}\le r_ir_j
 \le1+\sqrt{s/3}<\frac{25}{24}.
\]
Here $s/3\le1/600<1/576$.
With $\sigma=\sum_i(r_i-\overline r)^2$, the same pairwise variance
identities as in the five-point proof give
\[
 \sigma=\frac16\sum_{i<j}
       \frac{(t_i-t_j)^2}{(r_i+r_j)^2}<\frac{6s}{23},
 \qquad (r_i-r_j)^2<\frac{12s}{23}\le\frac3{1150}.
\]
The strengthened general distance bound gives $d_{ij}\ge12/25$.
Consequently,
\[
 \sin^2\frac{\theta_i-\theta_j}{2}
 >\frac{12/25-3/1150}{25/6}>\frac1{10};
 \qquad \left|\cot\frac{\theta_i-\theta_j}{2}\right|<3.
\]
Order the lifted angles and let $\theta^*$ be the equally spaced
six-point angles with the same mean. The circle logarithmic product
has Hessian at most $-3I$ on the mean-zero subspace. Its gradient
at the actual angles is given by \eqref{eq:angular-stationary}, with
six indices in place of five. The incidence matrix of $K_6$ has norm
$\sqrt6$, while, for $v_i=r_i-\overline r$,
\[
 \sum_{i<j}(v_i-v_j)^4
 =6\sum_i v_i^4+3\sigma^2\le8\sigma^2.
\]
Thus the angular error satisfies
\[
 \eta:=\norm{\theta-\theta^*}
 \le\frac4{\sqrt3}\frac{\sigma h}{d}
 <\frac73\frac{6s}{23}\,3\,\frac{25}{12}
 =\frac{175s}{46}\le\frac7{368}<\frac1{50}.
\]
Every half-angle error is at most $\eta/\sqrt2<1/70$, since
$\sqrt2>7/5$.
As sine is $1$-Lipschitz, the absolute half-angle sines in the three
cyclic separation classes are greater than $17/35$, $59/70$, and
$69/70$, respectively; here $\sqrt3/2>6/7$.
Multiplication by $4r_ir_j>23/6$ gives
\[
 d_{ij}>\frac{9}{10},\qquad d_{ij}>\frac83,\qquad
 d_{ij}>\frac{100}{27}
\]
in those three classes. Hence their Laplacian weights are less than
$a_0=10/9$, $b_0=3/8$, $c_0=27/100$.
The comparison circulant Laplacian has nonzero eigenvalues
$a_0+3b_0+2c_0$, $3a_0+3b_0$, and $4a_0+2c_0$, so
\[
 \lambda_{\max}(L)\le\frac{2243}{450}<5.
\]
This contradicts $L(t-\one)=5(t-\one)$ and $s>0$.
All numerical comparisons are rational inequalities or follow by
squaring positive rational numbers.
\end{proof}


\subsection{Why a six-point global argument needs a different target}
The inverse-distance Laplacian of the unit regular polygon similarly
has eigenvalues $k(n-k)/2$, $0\le k\le n-1$.
Indeed, its $k$th eigenvalue is
$\frac12\sum_{r=1}^{n-1}
|\sin(k\pi r/n)/\sin(\pi r/n)|^2$; expansion of
$|\sum_{j=0}^{k-1}e^{2\pi i jr/n}|^2$ and Fourier orthogonality
give $k(n-k)/2$.
The largest eigenvalue is $\frac12\lfloor n^2/4\rfloor$.
The strict spectral margin below $n-1$ holds for $2\le n\le6$, is zero
at $n=7$, and is absent for $n\ge8$. Thus the local exclusion by strict
Laplacian comparison cannot be copied unchanged for $n\ge7$.

For completeness, Chalk's competitor consists of the origin and a
regular $(n-1)$-gon of radius $\sqrt{n/(n-1)}$. It has
\begin{equation}\label{eq:Chalk-competitor}
 \Delta_n=(n-1)^{n-1}
       \left(\frac n{n-1}\right)^{n(n-1)/2}.
\end{equation}
At $n=6$ its ratio to $6^6$ is
$6^9/5^{10}=10077696/9765625>1$.
For real $x\ge6$, twice the logarithm of this ratio is
$g(x)=x(x-3)\log x-(x-1)(x-2)\log(x-1)$, whose derivative is
\[
 g'(x)=(2x-3)\log\frac{x}{x-1}-1
 \ge\frac{2x-5}{2x-1}>0.
\]
The lower bound uses
$\log(1+y)\ge2y/(2+y)$ for $y\ge0$, which follows by
differentiation. Hence \eqref{eq:Chalk-competitor} exceeds $n^n$
for every integer $n\ge6$.
Proposition~\ref{prop:six-gap} therefore separates the regular hexagon
from possible global maximizers, rather than proving that it is
globally optimal. A full six-point sharp bound would require a new
global comparison and additional constraints beyond the
two-dimensional projection.


\appendix
\section{Exact arithmetic and reproducibility}\label{app:certificate}
\subsection{Interval operations}
All endpoints in the program are instances of Python's
\texttt{fractions.Fraction}. Addition, subtraction, multiplication,
positive-part evaluation and squaring use their elementary outward
interval formulas. Division is performed only when the denominator
interval has positive lower endpoint.

For a nonnegative rational $v=a/b$, an outward square-root interval
with denominator $2^{40}$ is obtained from
\[
 m=\operatorname{isqrt}\left(\left\lfloor
              a\,2^{80}/b\right\rfloor\right).
\]
Its lower endpoint is $m/2^{40}$; the upper endpoint is this value
when it squares exactly to $v$, and otherwise $(m+1)/2^{40}$.
For a nondegenerate interval the lower and upper endpoint calculations
are applied separately. Thus no floating-point approximation enters
a mathematical decision. Wall-clock timings in the JSON output have
no role in certification.

The factor $20-12s-4q/5$ is nonnegative in the target domain.
The other factor is also positive: in A,
\[
 D_*=15+x^2(-180+5760x-15600x^2)\ge15-180/16>0,
\]
since $0\le x\le1/4$.
In B, put $T=5+s\in[25/4,20/3]$; then
$D_*=T(-12T^2+144T-405)/5>0$ by checking the endpoints of the
concave quadratic. Clipping interval factors to their nonnegative
parts is therefore valid, including rectangles extending outside the
curved domain boundary.

\subsection{Transcendental and algebraic constants}
The only transcendental checks concern logarithms.
The coefficient $c$ has the positive-series representation
\[
 c=\frac15-\frac4{25}\sum_{j=0}^{\infty}
                \frac1{(2j+1)5^j}.
\]
After $N$ terms the series tail is at most
$5/(4(2N+1)5^N)$. The program uses $N=24$.
For rational $x\ge1$, with $y=(x-1)/(x+1)$,
\[
 \log x=2\sum_{j=0}^{\infty}\frac{y^{2j+1}}{2j+1}.
\]
After $N$ terms, its tail is at most
$2y^{2N+1}/((2N+1)(1-y^2))$.
Using $N=120$ at $x=5$ and $x=41/6$ certifies
\eqref{eq:c-bounds} and
$H(R)-\log R-k(R^2-5R+5)^2>0$.

The checks in \eqref{eq:two-steps}--\eqref{eq:spectral-gap}
use the algebraic constants \eqref{eq:algebraic-trig}, each enclosed by
the integer-square-root procedure above.
No numerical value of $\pi$ or trigonometric series is required.
The underlying bounds $\sin x\le x$ and $\cos x\ge1-x^2/2$ follow
by integration of $\cos x\le1$ and $\sin x\le x$ for $x\ge0$.

\subsection{Polynomial expressions and source files}
To make the interval encoding explicit, in A the functions in
\eqref{eq:scalar-definitions} simplify to
\begin{equation}\label{eq:branch-A}
\begin{aligned}
 f(20x^2)&=x^2\bigl(-10+450C-500k
             +x(-300C+3000k-8500kx)\bigr),\\
 D_*(20x^2)&=15+x^2(-180+5760x-15600x^2).
\end{aligned}
\end{equation}
In B, with $T=5+s$,
\begin{equation}\label{eq:branch-B}
\begin{aligned}
 f(s)&=\left(-\frac12+\frac{15C}{2}\right)s
       +5C(T^2/5-5)
       -k(T^3/5+24T^2-275T+750),\\
 D_*(s)&=T(-12T^2+144T-405)/5.
\end{aligned}
\end{equation}

\noindent\begin{minipage}{\textwidth}
The ancillary files supplied with this manuscript are:
\begin{itemize}
\item \texttt{verify\_mordell.py};
\item \texttt{verification\_result.json};
\item \texttt{algebra\_identity\_check.py};
\item \texttt{algebra\_identity\_result.json};
\item \texttt{generalization\_check.py};
\item \texttt{generalization\_result.json}.
\end{itemize}
\end{minipage}
From their directory, the full certificate is reproduced by
\begin{verbatim}
python verify_mordell.py --output verification_result.json
python algebra_identity_check.py
python generalization_check.py
\end{verbatim}
The first command must be run with both regions, its default.
The program exits unsuccessfully if any scalar region is unresolved
or a constant comparison fails. Explicit exception checks remain active
under Python's optimization option \texttt{-O}.

The second script independently expands \eqref{eq:moment-expansion}
as formal polynomials over $\Q$, substitutes the mixed-moment
relations, and checks all four identities in
\eqref{eq:distance-moments}. It also derives
\eqref{eq:branch-A}--\eqref{eq:branch-B} from
\eqref{eq:scalar-definitions} and checks exact point encodings.
Its point checks are an encoding audit; the covering proof for region A
of \eqref{eq:certificate} is performed by the interval program, while
region B is handled by \eqref{eq:branch-b-analytic}.
The third script checks the general Gram coefficients and supporting
polynomial identities, the six-point rational constants, and the
exact six-point value in \eqref{eq:Chalk-competitor}.
The recursive identities and the spectral and variational arguments
are proved analytically above; finite script checks support those
proofs and do not replace them.

\section*{Acknowledgement of assistance}
AI assistance through OpenAI's Codex was used for mathematical
exploration, exact-arithmetic verification support, and manuscript
preparation. The author is responsible for the final mathematical
claims and for checking the accompanying proof and source code.

\begin{thebibliography}{9}
\bibitem{Mordell}
L.~J. Mordell, \emph{On a discriminant inequality},
Canadian Journal of Mathematics \textbf{12} (1960), 699--704.
\href{https://doi.org/10.4153/CJM-1960-063-5}
{doi:10.4153/CJM-1960-063-5}.


\bibitem{Chalk}
J.~H.~H. Chalk, \emph{A note on a discriminant inequality},
Canadian Mathematical Bulletin \textbf{3} (1960), no.~1, 7--9.
\href{https://doi.org/10.4153/CMB-1960-002-x}
{doi:10.4153/CMB-1960-002-x}.

\bibitem{Wang}
J.~Wang, \emph{The Five-Point Case of Mordell's Discriminant Inequality},
arXiv:2610.00358v1, 2026.
\url{https://arxiv.org/abs/2610.00358}.
\end{thebibliography}
\end{document}
